\section{A conditional route to quasi-polynomial recognition}
\label{sec:research-agenda}

The two results remove identifiable costs from a larger recognition
argument. Theorem~\ref{thm:raw-closure-equivalence} replaces every successful
raw singleton epoch by one source query and one polynomial compilation;
its cost need not contain a factor exponential in the number of raw
batches. The orbit index replaces repeated schedule discovery for a fixed
interval relation by one verified preparation. Neither result chooses the
next presentation or surface. A global complexity theorem must account
for that choice, for the number of objects chosen, and for their full
encoded descriptions.

\subsection{Three obligations that cannot be interchanged}

\paragraph{Deterministic discovery.}
A short certificate with a polynomial verifier need not be findable in
polynomial time. For the algebraic operation studied here, this gap is
closed: a prescribed survivor set is decided by Horn closure, and every
one-element survivor set can be tried in polynomial time. For normal
surfaces and hierarchy moves, a comparable discovery theorem is still
required. Counting or interpreting the components of a supplied normal
surface, as in AHT~\cite{aht}, does not find a suitable compressing disc or
prove that a small candidate list contains one. Candidate generation must
charge rejected candidates and unsuccessful subsidiary searches, including
those that never produce a certificate.

\paragraph{Hierarchy depth and the generated frontier.}
After each certified normalization or geometric cut, another local query
may be cheap, yet exponentially many distinct states may arise. One needs
a bound on the total states generated, or a depth bound together with a
branching bound, or a proved bound on the frontier after justified
identifications. Memoization provides a bound only after the number of
keys is bounded. Likewise, a polynomial bound on the size of each boundary
pattern does not bound the number of admissible patterns. The hierarchy
outline in the announced algorithm~\cite{lackenby-talk} and the detailed
2026 development~\cite{lackenby-2026} motivate this separation; the latter
does not supply all costs needed for the proposed implementation.

\paragraph{Bit complexity, provenance, and attachments.}
The cost of a state includes its reachable word grammar, binary lengths,
source-slot records, triangulation or handle data, verified trace, and
the maps specifying how boundary and attachment marks are identified.
Generator provenance and the evidence connecting the native input to the
original diagram are part of the state. So are historical masks, retained
arena nodes, and caches when operations can inspect them. Polynomial
compressed word and curve primitives~\cite{schleimer,lackenby-curves}
apply to the size of their supplied encodings. They cannot be charged
against a smaller crossing count without a bound relating the two.

The following accounting statement makes a sufficient target precise. It
is conditional on a complete recognition scheme, rather than a claim that
the scheme described in its hypotheses has been constructed.

\begin{proposition}[A sufficient global accounting contract]
\label{prop:qp-accounting}
Let $n\ge2$ be the bit length of a knot diagram encoding. Suppose a
deterministic recognition scheme has sound decomposition rules, complete
successor generation, and correct terminal tests. Suppose its whole
execution processes at most $F(n)$ state occurrences, counting repeated
states unless their suppression is justified. At each occurrence assume:
\begin{enumerate}
\item all discovery work, including rejected candidates, takes at most
      $D(n)$ bit operations;
\item all representations and relevant historical data have at most
      $P(n)$ bits, and their construction, verification, terminal tests,
      and composition cost $P(n)^{O(1)}$;
\item at most $Q(n)$ prepared local queries are made, each costing
      $P(n)^{O(1)}$ bit operations.
\end{enumerate}
If the initial faithful conversion costs $n^{O(1)}$, then, for a fixed
constant $c$, the running time is bounded by
\[
 n^{O(1)}+F(n)\bigl(D(n)+(1+Q(n))P(n)^c\bigr).
 \tag{A}\label{eq:global-accounting}
\]
In particular, quasi-polynomial bounds on $F,D,P,Q$ imply a
quasi-polynomial bound for this complete scheme.
\end{proposition}

\begin{proof}
Charge preparation and verification when each state is processed, then
charge its discovery and queries. Summing these bounds gives
\eqref{eq:global-accounting}. A finite product or fixed power of functions
bounded by $\exp((\log n)^{O(1)})$ has the same form. Completeness and
correct terminal interpretation are hypotheses; the estimate does not
derive them from certificate verification.
\end{proof}

For example, depth $h(n)$ and at most $b(n)$ generated successors per
state give $F(n)\le\sum_{i=0}^{h(n)}b(n)^i$. Polynomial branching and
$h(n)=O((\log n)^2)$ yield $\exp(O((\log n)^3))$ states. To obtain the
stronger $\exp(O((\log n)^2))$ by this particular argument would require
$h(n)=O(\log n)$, or a compensating smaller frontier. If at most $W(n)$
state occurrences are processed at each level, the alternative bound is
$F(n)\le(h(n)+1)W(n)$. Each formulation still needs the independent
$D,P,Q$ bounds. A complete fallback preserves correctness but contributes
its own running time whenever reached; the new local results do not
remove that contribution.

\subsection{Algebraic questions after saturation}

\paragraph{1. Which source presentations admit small survivor sets?}
For a fixed source define
\[
 \kappa(R)=\min\{|S|:S\subseteq X,\ \cl(S)=X\}.
\]
Enumeration gives a concrete quasi-polynomial procedure for the query
$\kappa(R)\le O(\log n)$ on polynomially encoded sources. The next study
should identify presentation classes with a proved small value, together
with a terminal method for their residual presentations. Reducing to a
small alphabet alone is not a recognition algorithm. General uniquely
restricted matching optimization has hardness results~\cite{golumbic};
whether useful parameterized structure survives the source-word
restrictions is a separate question. Any proposed improvement should
state its parameter and compare against the explicit
$\binom rk(r+m+I)$ enumeration bound.

This question must distinguish original meridians from transformed
generators. On standard unreduced crossing words, the implications recover
the familiar Wirtinger propagation mechanism~\cite{wirtinger}. When
crossing labels coincide, freely reducing those words can introduce
additional implications; ordinary diagrammatic coloring is then not
automatically identical to source closure. After Whitehead or other
verified presentation transformations, current generators also need not
be meridians. Corollary~4.3 of~\cite{incompatibility} supplies unknot
diagrams with arbitrarily large diagrammatic Wirtinger number. Thus a
fixed-diagram bounded-seed completeness conjecture for ordinary coloring
is already false. A useful new theorem must specify the permitted
transformations and charge their discovery.

\paragraph{2. Can cancellation boundaries be controlled?}
The example $xy,xy^2$ shows that a single free-reduction boundary can
enable eliminations absent from the original closure. A precise next
problem is to parameterize the interactions across such boundaries:
for instance by a stated bound on the number of participating source
factors or on an explicitly defined interaction graph. The desired
theorem would construct a finite successor list, bound its total encoded
size, and prove that some listed successor preserves a successful
continuation. Merely bounding the number of letters cancelled in an
expanded word would be unsuitable for highly compressed inputs.
Counterexample families should force cancellation across many shared
grammar subwords, so that a proposed parameter cannot hide an exponential
number of distinct contexts.

\paragraph{3. Which witness choices preserve later search behavior?}
Same-donor rematching preserves the induced map in the free group; closure
over all source rows may choose another donor set and need not preserve
that map. Can one optimize a closure witness for compiled size, attachment
description, or the next verified move without enumerating exponentially
many donor sets? A first falsifiable target is an efficiently computable
cost model for a restricted source class, with either a proved
approximation bound or a counterexample to greedy choice. Experiments
must compare semantic endpoints when donor sets agree, and complete
subsequent search outcomes when they differ. Equality of immediate ranks
does not establish interchangeability of these states. Classical matching
acyclicity and Horn propagation~\cite{kozen,dowling} remain the underlying
tools, rather than additional novelty claims.

\paragraph{4. Is there a potential controlling successive normalized epochs?}
The removed raw-phase factor should now be replaced by an explicit measure
of the transformations that remain. A promising theorem would associate
each normalized state with a potential and show a decrease, a bounded
number of neutral moves, or a balanced decomposition along every required
branch. The potential must account for encoded grammar complexity as well
as generator count: long sequences of rank-preserving transformations are
not ruled out by rank reduction. Record proposed potentials on adversarial
families and actively search for increasing cycles before using them in a
depth argument. A bound observed on successful examples is not a bound on
the failed branches required by deterministic recognition.

\subsection{Geometric representations and complete discovery}

\paragraph{5. Can the prepared point map be accelerated further?}
The present index stores the entire minimum set as at most $g$ ordinary
intervals and answers canonical minimum rank and selection in
$O(B\log(g+2))$ bit work. An arbitrary point query still follows the
contraction and truncation program; its stated bound is
$O(sB^2\log(k+2))$. These are different query problems. A concrete next
target is near-linear preparation in the stored trace size, followed by
polylogarithmic dependence on $s$ for point transport, with binary
arithmetic charged. Test whether balanced composition of partial affine,
reflection, and remainder maps produces a controlled representation, or
whether its number of pieces necessarily grows rapidly. The compactness
of the minimum set alone does not bound this composition. Existing
compressed tracing histories~\cite{erickson-nayyeri} are essential prior
art for this question.

\paragraph{6. Which attachment descriptions stay compact under cutting?}
Minimum representatives identify components of one fixed source relation.
They neither identify components across different triangulations nor
encode arbitrary permutations of attachment points. Choose an explicit
attachment language---for example, ordered marked intervals with specified
partial isometries and source charts---and prove closure of that language
under each proposed cut, gluing, and simplification. The measurable target
is a bound on chart count, map size, coordinate bits, and provenance length
after a sequence of operations. A counterexample with exponentially many
required chart pieces would refute that representation choice even if all
individual component queries remain polynomial. Signed variants must
transport the established geometric orientation data; an interval
reflection by itself is not an orientation character of a surface.

\paragraph{7. What makes a useful surface deterministically discoverable?}
The next hierarchy implementation should specify one exact move and its
input contract, then prove both a complete candidate description and a
producer with an encoded running-time bound. This is preferable to
assuming an unspecified oracle for a useful normal surface. Separate the
number of surface candidates, the weight bits of each, and the boundary
patterns created by cutting. AHT and the new index handle supplied
component data; the curve algorithms can handle supplied compressed curves
\cite{aht,lackenby-curves}. The missing theorem must explain why the
necessary supplied objects can be found within the budget. The 2026
hierarchy preprint~\cite{lackenby-2026} gives a serious source of structural
lemmas, but each imported statement must retain its hypotheses and cost
parameters.

\paragraph{8. Can equivalent marked frontiers be recognized and counted?}
A frontier compression theorem needs an equivalence relation preserving
all permitted continuations, a procedure to recognize that equivalence,
and a bound on the number of resulting classes. Equal component counts,
equal lists of minimum ranks, or isomorphic unmarked graphs need not
preserve attachment maps or peripheral information. Start with a restricted
surface class and a fixed attachment language from Question~6. Construct
pairs with identical coarse summaries but different marked behavior; these
are direct tests of any proposed state key. Only after such distinctions
are resolved can canonicalization justify a smaller $W(n)$ in the global
contract.

\paragraph{9. Can marked component data be returned without listing components?}
The most direct next geometric extension is to associate the maintained
weighted terminal runs with canonical minimum keys. Given one marked
point, return its component's Euler characteristic, boundary information,
or complete normal coordinate vector, using compressed weight transport.
The target is preparation polynomial in the source trace and requested
weight channels, followed by queries that depend on the encoded answer
size. The minimum-set interval bound does not automatically bound the
number of weight changes within those intervals. Prove an additional
bound on the joint representation, or exhibit a family in which the
chosen weight language requires too many pieces. The signed-cover
corollary supplies path parity, provided that the geometric character
has already been certified; it does not replace this weighted analysis.

\paragraph{10. Which proof obligations can be mechanically formalized first?}
A focused formalization should begin with the positive-pivot matching
injection, the same-donor free-group homomorphism argument, and the
orbit minimum-set invariant. Each has a small mathematical interface and
an independent finite oracle already available for examples. A second
stage should prove refinement from immutable grammar masks and the
queue invariant to the mathematical closure, and from the six verified
interval operations to the retained transport program. The distinction
between equality of raw strings and equality in a free group must appear
in the formal statement. No machine-checked proof is included in this
delivery: the supplied proofs are mathematical manuscripts, and the
independent certificate replayers and tests remain software evidence.

\subsection{An adversarial validation and integration programme}

The next experiments should test explicit claims and expose their limits.
They should retain the pinned baseline and independently verified outputs
used in this continuation~\cite{repo}. Three complementary designs would
make progress assessable.

\paragraph{Algebraic stress tests.}
Extend the finite matrix audit with structured unique-matching families,
bad pivot orders, repeated donor roots, and entries of large binary size.
For word-level tests vary signs independently of absolute incidence,
include empty slots and sparse generator names, and require same-donor
flattening to preserve every freely reduced generator image and residual
relator. Include presentations in which one normalization changes closure,
and knot families built from the Wirtinger obstruction examples. Record
closure work, compiled nodes, retained historical storage, normalization
boundaries, and all explored branches separately. The matrix checks
support implementation correctness; they do not replace the proof or
model the distribution of knot diagrams.

\paragraph{Orbit and attachment stress tests.}
Vary endpoint bit length independently of pairing count, trace length,
gap count, and query count. Combine reflections, periodic carriers, and
separated static gaps; compare different valid traces for identical
minimum maps and canonical order. Test huge component counts with few
minimum intervals, as well as families attaining the interval bound.
Cross-check small cases by explicit union--find and large structured cases
by independent formulas. For performance, compare the same requested
representatives and include a baseline sharing the original orbit count.
Report verification and preparation separately from prepared queries and
also as a complete workload, so the break-even query count is visible.

\paragraph{End-to-end and negative tests.}
Use held-out diagram families, verified diagram transformations, and
nontrivial knots, rather than only successful local certificates. Report
success, certified rejection, and resource exhaustion separately.
Mutations should target source bindings, donor reuse, ranks outside the
minimum set, attachment maps, incomplete traces, and arithmetic limits.
Record actual encoded sizes, peak retained memory, source hashes, raw
timings, and discovered-state counts. No finite timing fit establishes
the exponent in an asymptotic theorem; the useful outcome is either a
counterexample to a proposed bound or data identifying which obligation
in~\eqref{eq:global-accounting} dominates.

The resulting milestones are concrete: first, an exact and measured
normalization-boundary model; second, a closed language of marked
geometric states with composable provenance; third, a complete candidate
producer for a specified hierarchy move; and finally a proved frontier or
depth bound. The present saturation and indexing results supply reliable
local operations for that programme. The remaining global hypotheses are
research problems, not consequences of those local operations.
